\section*{Chapter \ref{ch:ll:cpp-for-latency-i-memory} --- C++ for Latency I: Memory}

\begin{solution}{exo:ll:cpp-for-latency-i-memory:1}
Offsets 0, 4, 8, 16 and 24; the data end at 26 and the size rounds up to 32: 16 bytes of padding for 16 of data. Sorted
(\texttt{double}, \texttt{int32}, \texttt{int16}, two \texttt{char}s): 16 bytes, no padding.
\end{solution}

\begin{solution}{exo:ll:cpp-for-latency-i-memory:2}
$50\,000 \times 64 = 3\,200\,000$ bytes, about 3.05~MiB: 782 pages of 4~KiB (two \omterm{def:ll:memory-hierarchy-and-caches:tlb}{huge pages}).
\end{solution}

\begin{solution}{exo:ll:cpp-for-latency-i-memory:3}
512~MiB is 131\,072 pages; at about \qty{1.3}{\micro\second} each, some 0.17 seconds at start-up, harmless. On the hot
path it is \qty{1.3}{\micro\second} added to whichever message first touches each page, 131\,072 times, at random moments.
\end{solution}

\begin{solution}{exo:ll:cpp-for-latency-i-memory:4}
\texttt{\textquotedbl{}AAPL\textquotedbl{}} fits the 15-character inline buffer: no allocation. The 18-character identifier allocates.
\texttt{reserve(0)} does not allocate. \texttt{v(1)} allocates storage for one element.
\end{solution}

\begin{solution}{exo:ll:cpp-for-latency-i-memory:5}
Replacing the \texttt{std::function} with a direct call or a template removes three (the lambda's copy of the fills and
the function's heap copy of the lambda with its capture); a fixed vector of fills removes three; a fixed table instead of
the map removes two. Together eight; the fixed identifier removes the last.
\end{solution}

\begin{solution}{exo:ll:cpp-for-latency-i-memory:6}
For the tail and for control. The heap's median hides its rare slow paths (consolidation, system calls, faults on new
pages); a pool has none, its slots are contiguous and pre-faulted, it refuses rather than grows, and its capacity is a
number the team chose and monitors.
\end{solution}

\begin{solution}{exo:ll:cpp-for-latency-i-memory:7}
Record the maximum of \texttt{live()} after each \texttt{create}. The churn benchmark's high-water mark is its 10\,000
live orders by construction. A sizing rule: take the year's measured peak, multiply by a safety factor for days worse
than any seen (1.5 to 2), round up to a power of two, and alarm when live objects pass a fraction (say 80\%) of it.
\end{solution}

\begin{solution}{exo:ll:cpp-for-latency-i-memory:8}
Both identifiers are short enough for the inline buffer, and the same message makes the same map key every time, so no new
node is ever created: the test cannot see the allocations of long identifiers and new keys. Test with production-shaped
messages (long identifiers, new keys, many fills) and assert the allocation count after warm-up.
\end{solution}

\begin{solution}{pb:ll:cpp-for-latency-i-memory:1}
\begin{enumerate}
\item $9 \times 50\,000 = 450\,000$ allocations a second.
\item The long identifier's string (1); the vector of fills growing to one, two and four elements (3); the map's new node
and its copy of the key (2); the lambda's copy of the fills (1); \texttt{std::function}'s heap storage for the lambda and
the copy of its captures (2).
\item The identifier's string and the map key's copy: two.
\item The vector's capacity doubles, so fills four to eight add one allocation, not four; beyond three the count grows
only logarithmically.
\item About \qty{25}{\nano\second} (a destroy and a create) on the committed measurement.
\item Nine pairs at about \qty{27}{\nano\second}: \qty{243}{\nano\second} a message, some 12 milliseconds a second, about
1\% of the thread.
\item About \qty{1.3}{\micro\second}; some three hundred faults would make \qty{400}{\micro\second}.
\item Consolidating fragmented free blocks; growing or trimming the heap with a system call; mapping a large block
directly, zeroed by the kernel; faulting in new pages.
\item A fixed string for the identifier, a fixed vector of fills, a fixed table of entries, a direct computation instead
of the callback.
\item Identifiers of 24 characters, eight fills, 64 entries; when exceeded, the handler refuses (the entry is not created)
and the condition must be counted and alarmed.
\item So that its pages are faulted in, and on the right node, before trading starts: by the thread that uses it, after
pinning (chapter~4).
\item 64 comparisons of short strings, some tens of nanoseconds; with thousands of entries it needs an open-addressing
hash table (chapter~19).
\item \textbf{Named result.} Nine allocations per message before, none after; on this laptop a \omterm{def:ll:cpp-for-latency-i-memory:fault}{page fault} costs about
\qty{1.3}{\micro\second}, forty times a write to a page already present.
\item It replaces \texttt{operator new} with a counter and asserts zero allocations per message after warm-up; a change
that allocates fails the build.
\item Interrupts and scheduling (chapter~13), cold caches after a quiet period, a lock held by another thread, a logging
call on the hot path.
\item Timestamp each stage of the slow message and correlate the slow ones with the allocator's statistics, the minor-fault
counter (\texttt{getrusage}) and the scheduler's context switches for the thread.
\item Refusing is a decision with a bounded cost and an alarm; allocating is an unbounded cost at an unknown moment.
\item In the test binary's counting global allocator (\texttt{firm\_arena::CountingAlloc}), and in the pool's API, which
returns \texttt{None} rather than growing.
\item About 57\,000 with a factor of 1.5, rounded up to 65\,536, with an alarm at 80\%.
\item Decide every byte the hot path will use before trading starts, touch it, and never ask the heap again.
\end{enumerate}
\end{solution}

\begin{solution}{iq:ll:cpp-for-latency-i-memory:1}
Padding shrinks, more objects fit per \omterm{def:ll:memory-hierarchy-and-caches:line}{cache line}, and the fields read together share a line: fewer misses per operation.
\iqlookfor{alignment rules and \omterm{def:ll:memory-hierarchy-and-caches:line}{cache lines}.}
\end{solution}

\begin{solution}{iq:ll:cpp-for-latency-i-memory:2}
The median hides the slow paths: consolidation, system calls to grow or trim the heap, zeroed pages from the kernel, faults
on new pages, and locks shared with other threads; their timing is not under the program's control.
\iqlookfor{tails and determinism rather than the average cost.}
\end{solution}

\begin{solution}{iq:ll:cpp-for-latency-i-memory:3}
An array of slots, each a union of the object's storage and a next pointer; a head pointer to the first free slot; create
pops and placement-constructs, destroy calls the destructor and pushes. When empty, return null and let the caller refuse
the work and raise an alarm; never fall back to the heap on the hot path.
\iqlookfor{an intrusive \omterm{def:ll:cpp-for-latency-i-memory:pool}{free list} and a deliberate exhaustion policy.}
\end{solution}

\begin{solution}{iq:ll:cpp-for-latency-i-memory:4}
Short strings are stored inside the object (15 characters in the GNU library), long ones on the heap, so allocation depends
on the data: tests with short identifiers never allocate, production with long ones does.
\iqlookfor{data-dependent allocation.}
\end{solution}

\begin{solution}{iq:ll:cpp-for-latency-i-memory:5}
A trap to the kernel on touching a page not yet backed by memory; it allocates and zeroes a page (microseconds). Allocate
everything at start-up, write every page (pre-fault), lock memory against paging (\texttt{mlockall}, chapter~13), and use
\omterm{def:ll:memory-hierarchy-and-caches:tlb}{huge pages}.
\iqlookfor{\omterm{def:ll:cpp-for-latency-i-memory:fault}{pre-faulting} and locking.}
\end{solution}

\begin{solution}{iq:ll:cpp-for-latency-i-memory:6}
A test build that replaces the global allocation functions with counting ones (or a counting global allocator in Rust),
drives the component with production-shaped data past its warm-up, and asserts that the count does not change; the test
runs in continuous integration on every change.
\iqlookfor{an automated assertion, with realistic data.}
\end{solution}
